% GENERATED FILE. Edit canonical lessons, then run npm run generate:books.
% canonical-source-hash: fnv1a64:ff4e2fffa2988a52
% canonical-lessons: BN-C493-desk, BN-C493-kerani, BN-C493-silmohor, BN-C493-pion, BN-C493-taip-kora

\chapter{At the Office}
\label{ch:bn-at-the-office}
\hlchaptermodality{493}

\begin{chapteropening}
\textbf{By the end of this chapter:} \emph{I can talk about desks, clerks, seals and typing in an office.}

Complete the last lesson of chapter 493: I can talk about desks, clerks, seals and typing in an office.
\end{chapteropening}

\section[\texorpdfstring{\bn{ডেস্ক} (ḍesk)}{ডেস্ক (ḍesk)}]{\bn{ডেস্ক} (ḍesk) — a desk}

\label{lesson:BN-C493-desk}



\begin{quote}
Before the new one: say the Bengali for a bell-metal smith, then the Bengali for to string together.
\end{quote}

\subsection*{You'll want to know: \bn{ডেস্ক}}
\textbf{\bn{ডেস্ক}} — \emph{ḍesk} — \textquotedblleft{}a desk\textquotedblright{}.

\begin{quote}
\bn{আমার} \bn{ডেস্কে} \bn{অনেক} \bn{ফাইল} \bn{আছে।}
\end{quote}

\begin{quote}
\emph{Āmār ḍeske ônek phāil āchhe.}
\end{quote}

\begin{quote}
There are lots of files on my desk.
\end{quote}

\begin{cousinweb}
\textbf{\bn{ডেস্ক}} (\emph{ḍesk}) is borrowed from English \emph{desk}.
\end{cousinweb}

\begin{grammarlens}[title={What you've built}]
One more word for describing what you lived and felt.
\end{grammarlens}

\subsection*{Your turn}
\begin{itemize}
\raggedright
  \item \emph{Say it:} \emph{ḍesk}
  \item \emph{Say it:} \emph{ḍesk}, once more
  \item \emph{Say it:} \emph{Āmār ḍeske ônek phāil āchhe.} — There are lots of files on my desk.
  \item \emph{Recall:} say the Bengali for a bell-metal smith, then the Bengali for to string together, then say \emph{ḍesk} again
\end{itemize}

\begin{tcolorbox}[breakable,colback=teal!4,colframe=teal!35!black,title={Before you move on}]
What does \bn{ডেস্ক} mean? (\textquotedblleft{}A desk\textquotedblright{}.) Say it once more.
\end{tcolorbox}
\section[\texorpdfstring{\bn{কেরানি} (kerāni)}{কেরানি (kerāni)}]{\bn{কেরানি} (kerāni) — a clerk}

\label{lesson:BN-C493-kerani}



\begin{quote}
Before the new one: say the Bengali for to string together, then the Bengali for a desk.
\end{quote}

\subsection*{You'll want to know: \bn{কেরানি}}
\textbf{\bn{কেরানি}} — \emph{kerāni} — \textquotedblleft{}a clerk\textquotedblright{}.

\begin{quote}
\bn{কেরানি} \bn{ফাইলটা} \bn{টেবিলে} \bn{রাখল।}
\end{quote}

\begin{quote}
\emph{Kerāni phāilṭā ṭebile rākhlô.}
\end{quote}

\begin{quote}
The clerk put the file on the table.
\end{quote}

\begin{cousinweb}
\textbf{\bn{কেরানি}} (\emph{kerāni}) is probably from Sanskrit \emph{karaṇa}, which Monier-Williams glosses as \textquotedblleft{}a writer, a scribe\textquotedblright{}; how the ending arose is not settled.
\end{cousinweb}

\begin{grammarlens}[title={What you've built}]
One more word for describing what you lived and felt.
\end{grammarlens}

\subsection*{Your turn}
\begin{itemize}
\raggedright
  \item \emph{Say it:} \emph{kerāni}
  \item \emph{Say it:} \emph{kerāni}, once more
  \item \emph{Say it:} \emph{Kerāni phāilṭā ṭebile rākhlô.} — The clerk put the file on the table.
  \item \emph{Recall:} say the Bengali for to string together, then the Bengali for a desk, then say \emph{kerāni} again
\end{itemize}

\begin{tcolorbox}[breakable,colback=teal!4,colframe=teal!35!black,title={Before you move on}]
What does \bn{কেরানি} mean? (\textquotedblleft{}A clerk\textquotedblright{}.) Say it once more.
\end{tcolorbox}
\section[\texorpdfstring{\bn{সিলমোহর} (silmohôr)}{সিলমোহর (silmohôr)}]{\bn{সিলমোহর} (silmohôr) — an official seal; a rubber stamp}

\label{lesson:BN-C493-silmohor}



\begin{quote}
Before the new one: say the Bengali for a desk, then the Bengali for a clerk.
\end{quote}

\subsection*{You'll want to know: \bn{সিলমোহর}}
\textbf{\bn{সিলমোহর}} — \emph{silmohôr} — \textquotedblleft{}an official seal; a rubber stamp\textquotedblright{}.

\begin{quote}
\bn{কেরানি} \bn{ফর্মে} \bn{সিলমোহর} \bn{লাগাল।}
\end{quote}

\begin{quote}
\emph{Kerāni phôrme silmohôr lāgālô.}
\end{quote}

\begin{quote}
The clerk stamped the form with the official seal.
\end{quote}

\begin{cousinweb}
\textbf{\bn{সিলমোহর}} (\emph{silmohôr}) probably joins English \emph{seal} to \bn{মোহর}, \textquotedblleft{}a seal\textquotedblright{}, usually traced to Persian \emph{muhr}: the same thing said twice, in two languages.
\end{cousinweb}

\begin{grammarlens}[title={What you've built}]
One more word for describing what you lived and felt.
\end{grammarlens}

\subsection*{Your turn}
\begin{itemize}
\raggedright
  \item \emph{Say it:} \emph{silmohôr}
  \item \emph{Say it:} \emph{silmohôr}, once more
  \item \emph{Say it:} \emph{Kerāni phôrme silmohôr lāgālô.} — The clerk stamped the form with the official seal.
  \item \emph{Recall:} say the Bengali for a desk, then the Bengali for a clerk, then say \emph{silmohôr} again
\end{itemize}

\begin{tcolorbox}[breakable,colback=teal!4,colframe=teal!35!black,title={Before you move on}]
What does \bn{সিলমোহর} mean? (\textquotedblleft{}An official seal; a rubber stamp\textquotedblright{}.) Say it once more.
\end{tcolorbox}
\section[\texorpdfstring{\bn{পিওন} (pion)}{পিওন (pion)}]{\bn{পিওন} (pion) — an office messenger; a peon}

\label{lesson:BN-C493-pion}



\begin{quote}
Before the new one: say the Bengali for a clerk, then the Bengali for an official seal.
\end{quote}

\subsection*{You'll want to know: \bn{পিওন}}
\textbf{\bn{পিওন}} — \emph{pion} — \textquotedblleft{}an office messenger; a peon\textquotedblright{}.

\begin{quote}
\bn{পিওন} \bn{ম্যানেজারের} \bn{জন্য} \bn{চা} \bn{আনল।}
\end{quote}

\begin{quote}
\emph{Pion mænejārer jônyô chā ānlô.}
\end{quote}

\begin{quote}
The office messenger brought tea for the manager.
\end{quote}

\begin{cousinweb}
\textbf{\bn{পিওন}} (\emph{pion}) is usually traced to Portuguese \emph{peão}, \textquotedblleft{}a foot soldier, a man on foot\textquotedblright{}.
\end{cousinweb}

\begin{grammarlens}[title={What you've built}]
One more word for describing what you lived and felt.
\end{grammarlens}

\subsection*{Your turn}
\begin{itemize}
\raggedright
  \item \emph{Say it:} \emph{pion}
  \item \emph{Say it:} \emph{pion}, once more
  \item \emph{Say it:} \emph{Pion mænejārer jônyô chā ānlô.} — The office messenger brought tea for the manager.
  \item \emph{Recall:} say the Bengali for a clerk, then the Bengali for an official seal, then say \emph{pion} again
\end{itemize}

\begin{tcolorbox}[breakable,colback=teal!4,colframe=teal!35!black,title={Before you move on}]
What does \bn{পিওন} mean? (\textquotedblleft{}An office messenger; a peon\textquotedblright{}.) Say it once more.
\end{tcolorbox}
\section[\texorpdfstring{\bn{টাইপ} \bn{করা} (ṭāip kôrā)}{টাইপ করা (ṭāip kôrā)}]{\bn{টাইপ} \bn{করা} (ṭāip kôrā) — to type}

\label{lesson:BN-C493-taip-kora}



\begin{quote}
Before the new one: say the Bengali for an official seal, then the Bengali for an office messenger.
\end{quote}

\subsection*{You'll want to know: \bn{টাইপ} \bn{করা}}
\textbf{\bn{টাইপ} \bn{করা}} — \emph{ṭāip kôrā} — \textquotedblleft{}to type\textquotedblright{}.

\begin{quote}
\bn{কেরানি} \bn{চিঠিটা} \bn{টাইপ} \bn{করল।}
\end{quote}

\begin{quote}
\emph{Kerāni chiṭhiṭā ṭāip kôrlô.}
\end{quote}

\begin{quote}
The clerk typed the letter.
\end{quote}

Past (he/she …): \bn{তিনি} \bn{ফর্মটা} \bn{টাইপ} \bn{করলেন}, \emph{tini phôrmṭā ṭāip kôrlen}.

\begin{cousinweb}
\textbf{\bn{টাইপ} \bn{করা}} (\emph{ṭāip kôrā}) joins \emph{type}, borrowed from English, to \bn{করা}, \textquotedblleft{}to do\textquotedblright{}.
\end{cousinweb}

\begin{grammarlens}[title={What you've built}]
One more word for describing what you lived and felt.
\end{grammarlens}

\subsection*{Your turn}
\begin{itemize}
\raggedright
  \item \emph{Say it:} \emph{ṭāip kôrā}
  \item \emph{Say it:} \emph{ṭāip kôrā}, once more
  \item \emph{Say it:} \emph{Kerāni chiṭhiṭā ṭāip kôrlô.} — The clerk typed the letter.
  \item \emph{Recall:} say the Bengali for an official seal, then the Bengali for an office messenger, then say \emph{ṭāip kôrā} again
\end{itemize}

\begin{tcolorbox}[breakable,colback=teal!4,colframe=teal!35!black,title={Before you move on}]
What does \bn{টাইপ} \bn{করা} mean? (\textquotedblleft{}To type\textquotedblright{}.) Say it once more.
\end{tcolorbox}
